\subsection{Seventh Letter}

\begin{quotex}
Here we find Rene Guenon admitting the existence of initiatic organizations in the West, whether derived from Masonry or the Christian Hermetism of the Middle Ages. He himself was “initiated” in a Western form. These organizations may not be known, but it is not out of the question for them to reveal certain things from time to time.

Obviously, Julius Evola was interested in the four men discussed (Eliphas Levi, Meyrink, Bo Yin Ra, and Kremmerz), all of whom Guenon dismisses. (Eckhart Tolle was inspired by Bo Yin Ra.) Even at the age of 50, after some 25 years of correspondence with Guenon, Evola does not understand some metaphysical teachings and seems to still be searching for some “initiation”.

Guenon is not taking into account that Hermetists are often tricksters, with a cultivated public persona that may be quite different from their real understanding. Esoteric writings are couched in symbols and often contain deliberate contradictions and hyperboles. This was often necessary, in fact, when there was danger of crossing political and religious authorities, to hide the true meaning from outsiders, while being understood by initiates.

Unfortunately, Guenon's obsession with initiation has in practice been deleterious since every Tom, Dick, and Harry believes he knows which “organization” is valid. Instead, they should follow Guenon's advice: learn the “theoretical knowledge” either from his writings or, even better, from the sources from which he himself got his theoretical knowledge. Once you have that knowledge, only then can you determine where to find its practical application. As he hinted, there are still initiatic organizations in the West based on Christian Hermetism; they just don't advertise on Craigslist for obvious reasons.

“Schuon's book” is The Transcendent Unity of Religion, which was published in French in 1948..

See Guenon/Evola Letter 7 Introduction for a more complete introduction. 

\end{quotex}
18 April 1949


Cairo, Egypt

About what you say in regards to Schuon's book, I don't see exactly how the affirmation of the metaphysical identity of the possible and the real could be an “error”, actually just the opposite. Nevertheless, if there is anyone to whom this should be attributed, it was I and not he, since, long before him, I devoted a whole chapter (the second) of the \emph{Multiple States of Being} to this problem.

As to the esoteric character of early Christianity, of which later Christianity was only an exteriorization (i.e., no longer having anything initiatic about it); we have no doubt about that, all the more since the Islamic tradition asserts it explicitly, claiming that Christianity, in its origins, was \emph{tariqa} [way] and not \emph{sharia} [law]. The absence of sharia is in fact evident from the moment that it had to supply it through an adaption of Roman law (whence “canon law” was derived), therefore with the contribution of something that was completely unrelated to Christianity (and it is necessary to note in this regard that the word in Arabic \emph{aqnun} is still used today, in contrast to sharia, to define every law that is not integrated in the tradition).

After my latest books (especially the Perspectives on Initiation and the Reign of Quantity, since in the Great Triad I only used two or three articles), there currently remain very few of my articles, as you noted, that have not yet been republished, at least among those that were intended to be copied in Ur. As for the articles on the \textbf{Fedeli d'Amore}, I must say that I had the intention for a long time to include them in a new edition of the Esoterism of Dante; I was not able to find the time to systematize it, but I did not abandon them and moreover it is likely that the edition is about to go out of print very quickly.

On the problem of Masonry, I believe that it will truly be very difficult for us to agree; but there are a few things that astonish me about what you say in this question. First of all, you make me say (without any qualification, for I had made quite clear that it concerned only the West) that “the only initiatic traditional organizations existing are the Compagnonnage\footnote{\url{https://en.wikipedia.org/wiki/Compagnons_du_Devoir}} and Masonry, and then you assert that I should not logically write that some Masons would seem to not take into account Oriental initiatic organizations, that otherwise exist and among which some have more or less numerous members in Europe itself. It stands to reason that it can at least write also for them. I add, to finish with this subject once and for all, that my writings can furnish to some people, unless they are associated with an initiatic organization, a theoretical knowledge of traditional doctrine that in itself is not unimportant. As to the Masons, they have had in this circle, in recent times, many more results than I myself had hoped.

Another thing: I said that in the Western world itself there still survive certain organization tied to Christian Hermetism and dating from the Middle Ages. If I have not stressed this more strongly, it is because they are so closed (one of those that I knew more in depth restricted its membership to just twelve) that the possibility of being admitted is in practice not even to be considered.

I come to another problem; if you take Masonry to be considered, or rather its origin, as an idea similar to what you express, I ask myself how you could have once had, as you had told me some time ago, the intention of developing a work on the rituals aimed at eliminating its anti-traditional elements that were introduced into it. Under these conditions it would be a totally useless task, and there is in that, I confess, something that is absolutely incomprehensible to me. [Before he was injured in Austria, Evola had been doing research on a book on Masonic rituals.]

However, what I would like you to take notice is this: the date of 1717 does not mark the origin of Masonry, but the beginning of its degeneration, something that is quite different. Furthermore, because we can speak of a utilization of “psychic residues” in this time, it would be necessary to suppose that operative Masonry had then ceased to exist, something not true, from the moment that it subsists still today in different countries, and that in England, between 1717 and 1813, it participated effectively to complete certain things and to straighten up others, at least in the measure in which that was still possible in a Masonry reduced to being only speculative. In reality, the schism of 1717 involved just four Lodges, while there still existed a number of much higher Lodges that did not take part in it. On the other hand, where a regular and continuous filiation exists, the degeneration did not interrupt the initiatic transmission; it only reduced its efficacy, at least in general lines, because in spite of everything there could always be exceptions. As to the anti-traditional action of which you speak, it would be necessary in this regard to make some precise distinctions, e.g., between the Anglo-Saxon and Latin Masons; but, in any case, that only proves the incomprehension of the great part of the members of one or another Masonic organization, simply a question of fact and not of principle. Fundamentally, what could be said is that Masonry was the victim of infiltrations of the modern spirit, as in the exoteric order even the Catholic Church is in its current state even to a greater degree.

Of course, it is that I do not want at all to attempt to persuade you or anybody, only to make you see that the problem is much more complex than what you seem to believe.

Regarding the “initiatic strains” you mention, without denying at all their existence (and what I just said of certain existent Hermetic groups could otherwise be connected to them), I have to say that, for many reasons, I am very skeptical in the great number of cases, and that even the examples you cited are among those that seem to me to be more than dubious. \textbf{Eliphas Levi} was a Mason, and certain English circles with which he was in relationship were, to summarize, study groups, without a proper initiatic character, and were reserved exclusively to Masons. Apart from that, I never found the least proof that he had received any other regular initiation, and all the stories that some have told about this can only be placed in the category of occultist fables.

\textbf{Meyrink} was certainly in the current of very numerous traditional ideas, especially of a Jewish source; but, not to mention that this presupposes precisely an initiation (I do not think that he was able to receive a truly Cabalistic initiation, the only thing without doubt to be taken in consideration in a case of the type), the burlesque manner and caricature with which he often presented these ideas leaves a truly sinister impression (a shame that I cannot tell you by letter all the efforts that I made to remediate certain detrimental consequences of The Green Face). Besides, his relations with the school of \textbf{Bo Yin Ra} (whose true origin I am probably the only one who knows, because I noticed that his own disciples did not know it) are not certainly a very favorable indication.

As for \textbf{Kremmerz}, there would be much to say, and it would require too much time. But what I saw of his writings and even certain rituals that were too clearly “fabricated”, gave me the impression of something of rather scant consistency, and that recalls the worst parts of Eliphas Levi's work; in any case, the different groups into which his disciples split give the impression of not knowing absolutely how to direct themselves, and some admit to finding themselves at a true impasse.

I will add that, every time that it concerned self-styled Egyptian organizations, there were more serious reasons to be suspicious of them, because nothing authentic is found there and not even, quite often, of any consistency. As for an initiation received outside of the ordinary way of joining a known organization, there are certainly some examples of them, but in these cases it was a question of extremely rare exceptions, and no one can be confident of finding himself in a similar condition to avoid a normal tie; to think of it differently would be self-delusion in a very serious way. As for me, since the age of 22 or 23, I have been attached with some initiatic organization both Oriental and Western, from which you can take account that the supposition you advanced could not in any way be applied to my situation.

The doctrinal question that you speak about at the end of your letter is, fundamentally, less difficult than what it seems at first view: every “true man” has instead realized all the possibilities of the human state, but each one following a way that is congenial to him and thanks to which he differentiates himself from the others. Moreover, if it were not so, how could there be a place here, in our world, also for other beings that have not reached this level? The same thing can also be applied, at another level, for the “transcendent man” or the \emph{jivan mukta}; but then it is a matter of the totality of the possibilities of all the states.

Only, what is real, as weird as it can seem, is in fact that beings who have reached the same level can then be, in a certain sense, “indistinguishable” from the outside, as far as it concerns corporeal appearance. It happens in effect that they encompass a “type” that no longer has any individuality, and that occurs above all for those who carry out certain special functions: the ”type” is then that of the same function, something can make one believe that it is always the same being to exercise it over the course of a period of several centuries, while the reality is something completely different.

\flrightit{Posted on 2012-07-30 by Cologero}

* * *

\begin{footnotesize}\begin{sffamily}

\texttt{Mihai on 2012-07-31 at 03:19 said: }

I intend to read a book by Meyrink- Der Engel vom westlichen Fenster. I am curious to draw some conclusions on the character, or at least his ideas. Evola certainly held him in high regard, but he is not the only one. Vasile Lovinescu, a Romanian follower of Guenon, who was also initiated on a Sufi path, mentions Der Engel vom… several times in his writings, calling it one time as “one of the few worthy esoteric novels of the 20th century”. 

I am wandering about what was the western initiation that Guenon received. I am only aware of the Sufi one.

Also, I am wandering if there are any such Hermetic organizations, that he mentioned, still existing in Europe today.


\hfill

\texttt{logres on 2012-07-31 at 06:27 said: }

The Simple Life of R. (P Charconac) claims Bo Yin Ra was the only European member of the Teshu-Maru, a Mongolian order.


\hfill

\texttt{Cologero on 2012-07-31 at 08:17 said: }

Mihai, have you read the Green Face? Guenon makes this strange statement: “a shame that I cannot tell you by letter all the efforts that I made to remediate certain evil consequences of The Green Face.” How bad can a novel be?

And it is not as though Guenon was casually dismissing Meyrink et al off the cuff; apparently he was deeply involved in investigating some of those figures.

Guenon was involved at that time of his life (22 or 23) with the French Hermetic school (e.g., Papus, Josephin Peladan, Fabre d'Olivet etc.). I assume that is the Western initiation he is referring to.

I believe that the major books of Valentin Tomberg and Boris Mouravieff are evidences of a Hermetic school; I have already pointed out some coincidences in their publication that cannot be coincidences. Tomberg traces back to an obscure Russian Hermetist (his name escapes me at the moment); Mouni Sadhu worked with the same teacher. Since the Rosicrucians escaped to the “East”, it should not be surprising that Russia maintained that tradition. The life and work of Joseph de Maistre in Russia attests to that.


\hfill

\texttt{Michael on 2012-07-31 at 11:08 said: }

I believe the Russian esotericist was G.O. Mebes.


\hfill

\texttt{Mihai on 2012-07-31 at 16:57 said: }

Cologero, I have not read any work by Meyrink until now. However, in the volume “Introduction to magic”, which contains some of the articles published by the Ur group there is an article called “The path of awakening according to Gustave Meyrink”- a collection of excerpts claimed to be taken from Der Golem and Das grüne Gesicht. In those excerpts there are clearly some equivocal expressions and attitudes mixed with, what I perceive to be, authentic esoteric thought. Of course, it is hard to pronounce a judgement basing oneself solely on a few excerpts. 

In Lovinescu's works (I did have not read all his works) I only came across praises and references on Der Engel vom westlichen Fenster. Perhaps this particular book is above the others, maybe Meyrink deepened his understanding over time (although I do not know the order in which his books appeared). I can only see when I will read some myself.


\hfill

\texttt{Avery Morrow on 2012-07-31 at 23:13 said: }

Bo Yin Ra is really an odd duck. His writing suggests a Theosophical connection; it's unclear what Guenon could have meant by his “true origin”.

\url{http://www.kober.com/name.php}

\url{http://en.wikipedia.org/wiki/Joseph\_Anton\_Schneiderfranken}


\hfill

\texttt{Mihai on 2012-08-01 at 04:45 said: }

I think he meant his country of origin or something.


\hfill

\texttt{logres on 2012-08-07 at 08:08 said: }

\url{http://www.alexandrededanann.net/boyinra\_fr.htm}

If anyone reads French, English still under construction – Bo Yin Ra and the “White Lodge” (which I think Guenon had a low view of).


\hfill

\texttt{Cologero on 2012-08-07 at 08:52 said: }

So the Dalai Lama was Bo Yin Ra's master whose portrait Guenon recognized? If so, what was Guenon's objection?


\hfill

\texttt{Cologero on 2012-08-08 at 13:09 said: }

Bo Yin Ra's master was Swami Narad Mani aka Hiran Singh. More later …


\hfill

\texttt{logres on 2012-08-09 at 17:54 said: }

Just got a Bo Yin Ra book today in the mail, so am anticipating your comments rather keenly. I must say that the excerpts of his on the official website don't strike me as typical Theosophical jargon.


\hfill

\texttt{Mr Zero on 2013-12-15 at 19:32 said: }

The society with only 12 members was called L'estoile internelle and was brought to Guenon's attention after he had left to Cairo in his correspondence with Reyor who was the editor of Etudes Traditionelles at the time. It is described in pages 80-82 of Sedgwick's book, they are previewable for free on google books. His sources are in the endnotes.


\hfill

\texttt{Cologero on 2020-08-01 at 18:04 said: }

My copy is in storage, so I assume the book referred to is “The Transcendent Unity of Religion”, which was published in French in 1949. Evola's confusion may be due to the word “possible” instead of “potential”. The Arabian commentators on Aristotle made that change, so that “potential existence” became the “possibility of being”, which Guenon and Schuon adapted. The latter way of expressing it seems to imply something contingent, i.e., it may or may not be; so how could that be equivalent to the real? On the other hand, Potential existence and actual existence are intimately related, when the potential is actualized. That makes it easier to see that potential is just as much being as is the actual.


\hfill
\end{sffamily}\end{footnotesize}
